% Part: intuitionistic-logic
% Chapter: tableaux
% Section: introduction

\documentclass[../../../include/open-logic-section]{subfiles}

\begin{document}

\olfileid[pt]{int}{tab}{int}

\olsection{Introdução}

!!^{tableau}s são certas árvores (que se ramificam para baixo) de !!{signed
  formula}s, isto é, pares formados por um sinal de valor de verdade ($\True$ ou
$\False$) e !!a{sentence}
\[
\sFmla{\True}{!A} \text{ ou } \sFmla{\False}{!A}.
\]
!!^a{tableau} começa com certo número de \emph{premissas}. Cada
!!{signed formula} adicional é gerada pela aplicação de uma das regras de
inferência. Algumas regras de inferência acrescentam uma ou mais !!{signed formula}s a uma
ponta da árvore; outras acrescentam duas novas pontas, produzindo dois ramos.
As regras produzem !!{signed formula}s nas quais a !!{formula} é
menos complexa do que a da !!{signed formula} à qual a regra foi
aplicada. Quando um ramo contém tanto $\sFmla{\True}{!A}$ quanto
$\sFmla{\False}{!A}$, dizemos que o ramo é \emph{fechado}. Se todos os
ramos de !!a{tableau} são fechados, o !!{tableau} inteiro é fechado. Um
!!{tableau} fechado constitui !!a{derivation} que mostra que o conjunto
de !!{signed formula}s usado para iniciar o !!{tableau} é
insatisfatível. Isso pode ser usado para definir uma relação $\Proves$:
$\Gamma \Proves !A$ se, e somente se, existe algum conjunto finito~$\Gamma_0 = \{!B_1,
\dots, !B_n\} \subseteq \Gamma$ tal que existe um
!!{tableau} fechado para as premissas
\[
\{\sFmla{\False}{!A}, \sFmla{\True}{!B_1}, \dots, \sFmla{\True}{!B_n}\}.
\]

Para a lógica intuicionista, precisamos tanto estender a noção de
!!{signed formula} quanto ajustar as regras dos conectivos. Além
de um sinal ($\True$ ou $\False$), as !!{formula}s nos
!!{tableau}s modais também têm \emph{prefixos}~$\sigma$. Os prefixos são
sequências não vazias de inteiros positivos, isto é, $\sigma \in
(\PosInt)^* \setminus \{\emptyseq\}$. Escrevemos tais prefixos
sem os $\tuple{\ }$ que os circundariam e separamos os
!!{element}s individuais por~$.$ em vez de $,$. Se $\sigma$ é um prefixo, então
$\sigma.n$ é $\sigma \concat \tuple{n}$; por exemplo, se $\sigma = 1.2.1$,
então $\sigma.3$ é $1.2.1.3$. Assim, por exemplo,
\[
\sFmla{\True}{!A \lif (!B \lif !C)}[1.2]
\]
é uma \emph{!!{signed formula} prefixada} (ou simplesmente uma \emph{
  !!{formula} prefixada}).

Intuitivamente, o prefixo nomeia um mundo de um modelo que poderia satisfazer
as !!{formula}s em um ramo de !!a{tableau}; se $\sigma$ nomeia
algum mundo, então $\sigma.n$ nomeia um mundo acessível a partir do mundo
nomeado por~$\sigma$.

Nos modelos intuicionistas, a relação de acessibilidade é reflexiva e
transitiva. Em termos de prefixos, isso significa que $\sigma$ é
acessível a partir do próprio $\sigma$, assim como qualquer prefixo que
estenda~$\sigma$, isto é, qualquer prefixo da forma
$\sigma.n_1.\cdots.n_k$. Introduzamos a notação $\sigma.*$ para
indicar o próprio $\sigma$ e qualquer extensão dele. Em outras palavras, os
prefixos $\sigma.*$ são precisamente os prefixos acessíveis
a partir de~$\sigma$.

\end{document}
